\chapter{Membuka Kotak Hitam: AI dan Visualisasi}
\label{ch:xai}

Kuliah AI and visualization saya ikuti pada musim dingin 2025/26, diberikan oleh kelompok Visual
Data Science di Linz. Kuliahnya bertumpu pada satu rumus yang muncul sejak pertemuan pertama: AI ditambah
visualisasi ditambah manusia sama dengan explainable AI. Penjelasan bukan sesuatu yang dikeluarkan
mesin begitu saja. Ia harus dirancang untuk seseorang, dan visualisasi adalah salah satu alat
terbaik untuk menyampaikannya.

Model yang paling akurat sering yang paling sulit dipahami. Untuk menyarankan film, itu tidak
masalah. Untuk menolak kredit, mendiagnosis penyakit, atau menyetujui desain sebuah reaktor, kita
perlu tahu alasan di balik keputusan. Bab ini membahas dua arah hubungan antara AI dan
visualisasi: visualisasi untuk memahami AI, dan AI untuk membuat visualisasi.

\section{Siapa yang butuh penjelasan, dan untuk apa}

Penjelasan yang dibutuhkan pengembang model berbeda dari yang dibutuhkan dokter, regulator, atau
orang yang permohonan kreditnya ditolak. Pengembang ingin menemukan bug dan bias. Pengguna ahli
ingin tahu kapan model bisa dipercaya. Orang yang terdampak ingin tahu apa yang harus diubah agar
keputusannya berbeda. Kuliah menyusun teknik-teknik XAI menurut pertanyaan pengguna: mengapa
prediksi ini, mengapa bukan yang lain, bagaimana kalau input diubah, dan contoh mana yang mirip.

Ada juga pertukaran yang sering digambarkan antara kinerja dan keterjelasan. Regresi linear dan
pohon keputusan kecil mudah dibaca, jaringan dalam tidak. \citet{rudin2019} berpendapat bahwa untuk
keputusan berisiko tinggi, kita sebaiknya memakai model yang bisa ditafsirkan sejak awal daripada
menjelaskan kotak hitam setelah jadi, karena penjelasan sesudahnya bisa saja tidak setia pada apa
yang sebenarnya dilakukan model.

Kuliah membedakan dua arah. \istilah{Vis for ML} memakai visualisasi untuk memahami, men-debug, dan
membandingkan model. \istilah{ML for Vis} memakai machine learning untuk membantu membuat visualisasi,
misalnya merekomendasikan jenis grafik atau menghasilkan dashboard dari permintaan dalam bahasa
biasa.

\section{Dasar-dasar visualisasi}

Visualisasi bekerja karena sistem visual manusia memproses pola dengan cepat. Tetapi tidak
semua saluran visual sama kuatnya. \citet{munzner2014} menyusun peringkat untuk data kuantitatif:
posisi pada sumbu yang sama paling akurat dibaca, disusul panjang, sudut, luas, lalu kecerahan dan
saturasi warna. Itulah sebabnya grafik batang biasanya lebih jujur daripada diagram lingkaran, dan
mengapa grafik tiga dimensi yang memiringkan sumbu sering menyesatkan. Prinsip Gestalt menjelaskan
bagaimana mata mengelompokkan objek: berdekatan, mirip warna atau bentuk, terhubung garis, atau
berada di dalam wadah yang sama.

Visualisasi juga bisa bercerita. Cerita yang digerakkan penulis membawa pembaca melalui urutan yang
sudah ditentukan, seperti artikel berita. Cerita yang digerakkan pembaca membiarkan orang menjelajah
sendiri. Banyak visualisasi terbaik memakai bentuk campuran: dimulai dengan narasi, lalu membuka
ruang untuk dijelajahi.

Kuliah memakai visualisasi untuk menjelaskan algoritme. Algoritme pengurutan bisa digambar sebagai
garis-garis yang bertukar tempat dari langkah ke langkah. $k$-means bisa dianimasikan sebagai pusat
yang bergeser. Teorema Bayes bisa dijelaskan dengan kotak berisi seratus orang. Contoh yang dibahas
adalah tes penyakit dengan sensitivitas 90 persen, tingkat positif palsu 9 persen, dan prevalensi 1
persen. Dari seribu orang, sepuluh sakit dan sembilan di antaranya terdeteksi, sedangkan dari 990
orang sehat, sekitar 89 mendapat hasil positif palsu. Peluang benar-benar sakit setelah hasil positif
hanya $9/(9+89)\approx9$ persen. Banyak orang, termasuk dokter, menebak jauh lebih tinggi. Gambar
seratus kotak kecil membuat hasil ini langsung terasa masuk akal.

\section{Menjelaskan lewat proyeksi}

Data berdimensi tinggi, termasuk aktivasi di dalam jaringan saraf, bisa diproyeksikan ke dua dimensi
untuk dilihat. PCA (\cref{ch:unsup}) mempertahankan varians global, t-SNE \citep{vandermaaten2008} mempertahankan tetangga
lokal, dan UMAP \citep{mcinnes2018} berusaha menjaga keduanya. Kuliah meluangkan waktu khusus untuk
salah baca yang umum pada t-SNE: hyperparameter seperti perplexity amat menentukan, ukuran kelompok
di peta tidak berarti apa-apa, jarak antar kelompok bisa menipu, noise acak kadang tampak berpola,
dan untuk melihat topologi kita mungkin butuh beberapa peta dengan setelan berbeda.

Salah baca itu bisa dipahami dari cara t-SNE bekerja. Untuk setiap titik di ruang asal, t-SNE mengubah jarak ke
titik-titik lain menjadi peluang bertetangga $p_{ij}$ dengan kernel Gauss, yang lebarnya dipilih per titik agar setiap
titik punya kira-kira sejumlah tetangga efektif yang sama, yaitu perplexity. Di peta dua dimensi, peluang $q_{ij}$
dihitung dengan distribusi Student-$t$ berekor tebal, lalu posisi titik di peta digeser untuk meminimalkan
$\KL(P\|Q)=\sum p_{ij}\log(p_{ij}/q_{ij})$. Karena lebar kernel disesuaikan per titik, kelompok yang rapat dan kelompok
yang renggang di ruang asal sama-sama tampil dengan ukuran serupa di peta. Dan karena KL ini menghukum keras tetangga
dekat yang dipisahkan, tetapi hampir tidak peduli pada titik jauh yang didekatkan, jarak antar kelompok di peta tidak
punya arti yang bisa diandalkan.

Proyeksi juga bisa menunjukkan proses yang berjalan dalam waktu. Kalau setiap keadaan sebuah proses,
misalnya setiap langkah algoritme, setiap posisi permainan catur, atau setiap epoch pelatihan
jaringan, diproyeksikan sebagai satu titik, lalu titik-titik itu dihubungkan, kita mendapat lintasan.
Lintasan pelatihan dua jaringan bisa dibandingkan: apakah keduanya menempuh jalan yang mirip, atau
satu berputar-putar sebelum menemukan solusi. Kelompok visualisasi di Linz juga mengembangkan alat
seperti ini untuk kimia, sehingga ahli kimia bisa menjelajahi ruang molekul yang dilihat oleh model
prediksi mereka.

\section{Menjelaskan dalam ruang input}

Sekelompok teknik menjelaskan model dengan melihat bagaimana keluarannya berubah terhadap input.
\istilah{Partial dependence plot} menampilkan efek rata-rata satu fitur terhadap prediksi, dengan
merata-ratakan atas nilai fitur-fitur lain. \istilah{Individual conditional expectation} menampilkan
kurva yang sama untuk setiap contoh secara terpisah, sehingga interaksi antar fitur yang tersembunyi
di balik rata-rata menjadi terlihat. Ambil model $f(x_1,x_2)=x_1x_2$ dengan $x_2$ bernilai $+1$ atau $-1$ sama seringnya.
Separuh kurva ICE untuk $x_1$ naik, separuh lagi turun, dan PDP, rata-ratanya, datar sempurna (\cref{fig:ice}). Kalau
hanya melihat PDP, kita akan menyimpulkan bahwa $x_1$ tidak berpengaruh sama sekali, padahal ia berpengaruh, hanya
arahnya bergantung pada $x_2$.

\begin{figure}[ht]
\centering
\begin{tikzpicture}
\begin{axis}[width=0.6\linewidth,height=4.4cm,axis lines=left,xmin=-1,xmax=1,ymin=-1.15,ymax=1.15,
  xlabel={\small $x_1$},ylabel={\small prediksi},xtick={-1,0,1},ytick={-1,0,1}]
  \foreach \b in {0.6,0.8,1.0}{
    \addplot[biru!60,domain=-1:1] {\b*x};
    \addplot[oranye!70,domain=-1:1] {-\b*x};
  }
  \addplot[black,very thick,dashed,domain=-1:1] {0};
  \node[label] at (axis cs:0.55,-0.3) {PDP};
\end{axis}
\end{tikzpicture}
\caption{Kurva ICE (garis tipis) untuk model dengan interaksi. Kurva naik untuk sebagian contoh dan turun untuk
sebagian lain, sehingga rata-ratanya, PDP (garis putus-putus), datar dan menyembunyikan pengaruh fitur itu.}
\label{fig:ice}
\end{figure}

LIME \citep{ribeiro2016} menjelaskan satu prediksi dengan membangun model sederhana yang hanya berlaku
di sekitar contoh itu. Contoh diganggu sedikit demi sedikit, misalnya dengan menyembunyikan
potongan-potongan gambar (\istilah{superpixel}), model asli ditanya untuk setiap versi, lalu model
linear dilatih untuk meniru jawabannya secara lokal. Secara formal, LIME mencari
\begin{equation}
g^*=\argmin_{g\in G}\ \sum_{z}\pi_x(z)\big(f(z)-g(z)\big)^2+\Omega(g),
\end{equation}
dengan $z$ contoh-contoh yang diganggu, $\pi_x(z)$ bobot yang makin kecil untuk $z$ yang makin jauh dari $x$, dan
$\Omega(g)$ penalti kerumitan, misalnya batas jumlah fitur yang boleh dipakai. Ini regresi linear berbobot dengan
regularisasi dari \cref{ch:data}. Bobot model linear itu menjadi penjelasan. Kelemahannya juga terlihat dari rumus
ini: hasilnya bergantung pada cara mengganggu contoh dan pada lebar $\pi_x$, sehingga dua setelan yang sama-sama masuk
akal bisa memberi penjelasan yang berbeda. Model
surrogate global melakukan hal serupa untuk seluruh data, dengan risiko bahwa model sederhana itu
tidak cukup setia.

Nilai Shapley berasal dari teori permainan kooperatif. Pertanyaannya: kalau sekelompok pemain bersama
menghasilkan keuntungan, berapa bagian yang adil untuk setiap pemain? Jawaban Shapley adalah
rata-rata kontribusi marginal pemain itu atas semua urutan kedatangan pemain. Dalam XAI, pemainnya
adalah fitur dan keuntungannya adalah prediksi. Ambil contoh dua fitur, A dan B. Tanpa fitur apa pun,
prediksi model 0. Dengan A saja 10, dengan B saja 20, dan dengan keduanya 40. Kalau A datang lebih
dulu, kontribusinya 10, dan kalau datang setelah B, kontribusinya 20. Rata-ratanya 15. Dengan
cara yang sama, B mendapat $(20+30)/2=25$. Jumlah keduanya tepat 40, prediksi penuh. Secara umum, untuk $n$ fitur dan fungsi nilai $v(S)$ yang memberi prediksi dengan hanya
memakai himpunan fitur $S$,
\begin{equation}
\phi_i=\sum_{S\subseteq\{1,\dots,n\}\setminus\{i\}}\frac{|S|!\,(n-|S|-1)!}{n!}\big(v(S\cup\{i\})-v(S)\big).
\end{equation}
Bobotnya adalah peluang bahwa, dalam urutan kedatangan yang acak, tepat fitur-fitur di $S$ datang sebelum $i$: ada
$|S|!$ cara menyusun yang datang lebih dulu, $(n-|S|-1)!$ cara menyusun yang datang belakangan, dari $n!$ urutan. Nilai
Shapley adalah satu-satunya pembagian yang memenuhi empat sifat sekaligus: jumlahnya sama dengan prediksi penuh
dikurangi prediksi kosong, fitur yang berperan identik mendapat bagian sama, fitur yang tidak pernah berpengaruh
mendapat nol, dan untuk jumlah dua model, nilainya juga dijumlahkan. SHAP \citep{lundberg2017} membuat perhitungan ini
praktis untuk model dengan banyak fitur, karena jumlah urutan yang mungkin tumbuh secara faktorial. Yang sering
terlupa adalah bahwa $v(S)$ harus didefinisikan: apa artinya model ``tanpa'' sebuah fitur? Biasanya fitur yang hilang
diganti nilai dari data latar, dan kalau fitur-fitur saling berkorelasi, penggantian itu bisa menciptakan contoh yang
tidak realistis, seperti orang berusia sepuluh tahun dengan gelar doktor.

Pertanyaan ``bagaimana kalau'' dijawab oleh \istilah{counterfactual explanation} \citep{wachter2017}: cari contoh $x'$
yang sedekat mungkin dengan $x$ tetapi menghasilkan keputusan yang diinginkan, misalnya dengan meminimalkan
$\lambda\big(f(x')-y'\big)^2+d(x,x')$. Jawabannya berbentuk kalimat yang langsung bisa dipakai: ``kredit Anda akan
disetujui kalau pendapatan Anda 400 euro lebih tinggi per bulan''. Penjelasan seperti ini tidak membuka isi model sama
sekali, sehingga cocok untuk hak atas penjelasan dalam GDPR (\cref{ch:manusia}). Tetapi ia harus dibatasi pada
perubahan yang mungkin: menyarankan seseorang menjadi lebih muda bukan penjelasan yang berguna.

\section{Menjelaskan lewat gradien}

Untuk jaringan saraf, gradien keluaran terhadap input memberi peta kepekaan langsung: piksel mana yang
paling mengubah skor kelas kalau digeser sedikit. \istilah{Saliency map} vanila \citep{simonyan2014}
menampilkan besar gradien ini. Hasilnya sering berbintik dan berisik. Integrated gradients
\citep{sundararajan2017} menjumlahkan gradien sepanjang jalur dari gambar acuan, misalnya gambar hitam,
ke gambar asli, dan memenuhi sifat bahwa jumlah atribusinya sama dengan selisih keluaran:
\begin{equation}
\mathrm{IG}_i(x)=(x_i-x_i')\int_0^1\frac{\partial f}{\partial x_i}\big(x'+\alpha(x-x')\big)\dd\alpha .
\end{equation}
Contoh kecil memperlihatkan mengapa sifat ini penting. Untuk $f(x_1,x_2)=x_1x_2$ di $x=(2,3)$ dengan acuan nol, di sepanjang
jalur $\partial f/\partial x_1=3\alpha$, sehingga $\mathrm{IG}_1=2\int_0^13\alpha\dd\alpha=3$, dan dengan cara yang sama
$\mathrm{IG}_2=3$. Jumlahnya 6, tepat $f(x)-f(0)$. Atribusi gradien kali input, yang memakai gradien di titik akhir saja,
memberi $3\cdot2=6$ dan $2\cdot3=6$, dengan jumlah 12, dua kali lipat dari keluaran yang ingin dijelaskan. Grad-CAM
\citep{selvaraju2017} memakai gradien terhadap peta fitur lapisan konvolusi terakhir untuk menimbang
peta-peta itu. Bobot peta ke-$k$ adalah rata-rata gradien skor kelas $y^c$ terhadap peta itu,
$\alpha_k=\tfrac1Z\sum_{i,j}\partial y^c/\partial A^k_{ij}$, dan peta panasnya $\mathrm{ReLU}\big(\sum_k\alpha_kA^k\big)$, sehingga menghasilkan peta panas kasar yang menunjukkan daerah gambar yang dipakai
model. Layer-wise relevance propagation \citep{bach2015} membagi skor keluaran mundur ke input
lapisan demi lapisan dengan aturan kekekalan relevansi.

Kuliah memakai contoh dari sebuah studi klasifikasi jamur. Peta Grad-CAM kadang menunjukkan bahwa model
melihat bagian yang memang penting secara biologis, seperti tudung dan bilah, tetapi kadang juga
latar belakang, misalnya rumput atau tangan yang memegang jamur. Peta seperti ini tidak membuktikan
bahwa model benar, tetapi amat berguna untuk menemukan jalan pintas yang tidak diinginkan.

Sebaliknya, peta yang tampak masuk akal juga belum tentu setia pada model. \citet{adebayo2018} mengusulkan uji kewarasan
yang sederhana: acak bobot jaringan lapisan demi lapisan, atau latih ulang dengan label acak, lalu lihat apakah peta
penjelasannya ikut berubah. Beberapa metode populer menghasilkan peta yang hampir sama untuk jaringan terlatih dan
jaringan acak. Metode seperti itu lebih banyak memperlihatkan tepi-tepi gambar daripada apa yang dipelajari model,
seperti detektor tepi yang kebetulan rapi. Pelajarannya berlaku untuk semua XAI: penjelasan juga harus diuji, bukan
hanya dinikmati.

\section{Menjelaskan lewat contoh dan konsep}

Manusia sering menjelaskan dengan contoh: ``ini mirip kasus yang kemarin''. Model pun bisa
menjelaskan prediksinya dengan menunjukkan contoh data latih yang paling mirip di ruang fiturnya.
\istilah{Prototype} adalah contoh yang mewakili sebuah kelompok, dan \istilah{criticism} adalah contoh yang
tidak terwakili dengan baik oleh prototype mana pun. Keduanya bersama memberi gambaran lebih jujur
tentang data.

Penjelasan berbasis konsep bertanya apakah sebuah konsep yang bermakna bagi manusia, misalnya
``bergaris'' atau ``berlatar salju'', tersimpan di dalam jaringan. \istilah{Network dissection} mencocokkan
unit-unit di dalam jaringan dengan konsep yang dianotasi. TCAV \citep{kim2018} mempelajari arah di
ruang aktivasi yang mewakili sebuah konsep dari contoh-contohnya, lalu mengukur seberapa peka
prediksi sebuah kelas terhadap arah itu. Jenis temuan yang dicari mirip contoh terkenal dari
\citet{ribeiro2016}, di mana penjelasan LIME mengungkap bahwa sebuah pengklasifikasi serigala dan
anjing husky sebenarnya terutama melihat salju di latar belakang.

\section{Visualisasi di era AI generatif}

Pertemuan terakhir membahas arah sebaliknya: AI yang membuat visualisasi. Model bahasa bisa menulis
kode grafik dari permintaan dalam bahasa biasa, menyusun dashboard, atau membuat penjelasan bagi
pengguna baru. Kuliah menyoroti risikonya. Gambar yang dihasilkan bisa tampak meyakinkan tetapi
salah secara data. Model membuat asumsi tentang maksud pengguna yang tidak pernah diucapkan. Dan
jawabannya bergantung pada gambaran yang dibentuk model tentang siapa penggunanya. Kuliah
mengusulkan antarmuka yang membuat asumsi-asumsi itu terlihat dan bisa dikoreksi pengguna.

\section{Di luar kelas}

Keterjelasan penting di setiap bidang yang dibahas buku ini, dengan alasan yang berbeda-beda.
Dalam kedokteran, peta Grad-CAM pada citra rontgen bisa mengungkap model yang melihat penanda
rumah sakit di sudut gambar alih-alih paru-paru. Dalam penemuan obat, atribusi pada atom-atom molekul
membantu ahli kimia memutuskan bagian mana yang perlu diubah. Dalam pemrosesan sinyal dan sistem
tertanam, kita perlu tahu sensor mana yang paling menentukan sebuah keputusan. Dalam simulasi,
keterjelasan bisa diambil sejauh-jauhnya: metode seperti SINDy (\cref{ch:penemuan}) tidak menjelaskan
model kotak hitam, melainkan menghasilkan persamaan yang bisa dibaca langsung. Itu barangkali bentuk
penjelasan yang paling disukai \citet{rudin2019}.

\section{Perkembangan riset}

Untuk model bahasa besar, riset keterjelasan bergeser ke \istilah{mechanistic interpretability}: mencoba membaca
mekanisme di dalam jaringan, bukan hanya atribusi pada input. Masalah utamanya adalah \istilah{polysemanticity}. Satu
neuron sering aktif untuk banyak konsep yang tidak berhubungan, karena jaringan diduga menyimpan lebih banyak fitur
daripada jumlah neuronnya, sebagai arah-arah yang saling bertumpuk di ruang aktivasi (\istilah{superposition}).
\citet{cunningham2023} melatih sparse autoencoder untuk merekonstruksi aktivasi internal model bahasa dari kombinasi
jarang arah-arah yang jauh lebih banyak dari jumlah neuron. Fitur yang ditemukan lebih monosemantik, satu fitur untuk
satu konsep, dan bisa dipakai untuk menunjuk fitur mana yang secara kausal bertanggung jawab atas sebuah perilaku.
Gagasannya adalah analisis faktor jarang dari \cref{ch:unsup}, diterapkan pada aktivasi jaringan.

Arah lain membangun keterjelasan ke dalam arsitektur, sejalan dengan saran \citet{rudin2019}. Concept bottleneck model
\citep{koh2020} memaksa jaringan memprediksi dulu konsep-konsep yang bermakna bagi manusia, misalnya ``ada taji tulang''
pada citra rontgen lutut, lalu memprediksi label akhir hanya dari konsep itu. Dokter bisa melihat konsep mana yang
dipercaya model, dan bisa mengoreksinya saat model dipakai. Koreksi seperti itu dilaporkan meningkatkan akurasi secara
nyata.

\sumberbab{Bab ini mengikuti enam kuliah AI and Visualization beserta labnya: dasar visualisasi,
menjelaskan algoritme, proyeksi, teknik penjelasan visual (PDP, ICE, LIME, Shapley, saliency, Grad-CAM,
contoh, konsep), dan AI generatif untuk visualisasi. Bacaan pendamping: \citet{molnar2022} untuk
interpretable machine learning dan \citet{munzner2014} untuk visualisasi.}
